%! The Secant, Cosecant, and Cotangent Functions — Unit 3, Lesson 3.11 — Homework
\documentclass[10pt]{article}
\usepackage{precalculus-article}
\usepackage{precalculus-boxes}
\newcommand{\TallMath}[1]{$\displaystyle #1\rule[-1.4em]{0pt}{3.2em}$}

\begin{document}

\pageheader{Unit 3, Lesson 3.11}{Homework}
\namedateperiod

\vspace{0.10in}

\begin{notesbox}{Problems 1--5}

\textbf{1.}\quad Evaluate each reciprocal trigonometric function using unit-circle values.

\vspace{4pt}
\renewcommand{\arraystretch}{1.5}
\begin{tabular}{|l|c|}
\hline
\textbf{Expression} & \textbf{Value} \\
\hline
$\sec\!\left(\dfrac{\pi}{6}\right)$ (use $\cos(\pi/6) = \sqrt{3}/2$) & \\
\hline
$\csc\!\left(\dfrac{\pi}{4}\right)$ (use $\sin(\pi/4) = \sqrt{2}/2$) & \\
\hline
$\cot\!\left(\dfrac{\pi}{3}\right)$ (use $\sin(\pi/3) = \sqrt{3}/2$, $\cos(\pi/3) = 1/2$) & \\
\hline
$\sec(\pi)$ (use $\cos(\pi) = -1$) & \\
\hline
$\csc\!\left(\dfrac{3\pi}{2}\right)$ (use $\sin(3\pi/2) = -1$) & \\
\hline
\end{tabular}

\vspace{0.12in}
\textbf{2.}\quad For each function, give the general formula for the vertical asymptotes
and the range.

\vspace{4pt}
(a)\quad $f(\theta) = \sec\theta$: \quad Asymptotes: $\theta = $ \underline{\hspace{3.5cm}}
\quad Range: \underline{\hspace{4cm}}

\vspace{6pt}
(b)\quad $g(\theta) = \csc\theta$: \quad Asymptotes: $\theta = $ \underline{\hspace{3.5cm}}
\quad Range: \underline{\hspace{4cm}}

\vspace{6pt}
(c)\quad $h(\theta) = \cot\theta$: \quad Asymptotes: $\theta = $ \underline{\hspace{3.5cm}}
\quad Range: \underline{\hspace{4cm}}

\vspace{0.12in}
\textbf{3.}\quad Decide whether each value is in the domain. Write ``Yes'' or ``No''
and explain in one sentence.

\vspace{4pt}
(a)\quad $\theta = \pi/2$ for $\sec\theta$: \underline{\hspace{1.5cm}}
Reason: \writeline

\vspace{4pt}
(b)\quad $\theta = 2\pi/3$ for $\csc\theta$: \underline{\hspace{1.5cm}}
Reason: \writeline

\vspace{4pt}
(c)\quad $\theta = \pi$ for $\cot\theta$: \underline{\hspace{1.5cm}}
Reason: \writeline

\vspace{0.12in}
\textbf{4.}\quad A graph of $y = \cos\theta$ is shown below with two zeros marked
at $\theta = \pi/2$ and $\theta = 3\pi/2$.

\vspace{4pt}
(a)\quad Near $\theta = \pi/2$, the value of $\cos\theta$ is very close to 0.
What happens to $\sec\theta = 1/\cos\theta$ as $\theta$ approaches $\pi/2$?

\writeline

\vspace{4pt}
(b)\quad Between $\theta = 0$ and $\theta = \pi/2$, $\cos\theta$ decreases from
1 to 0. Does $\sec\theta$ increase or decrease over this interval?

\writeline

\vspace{4pt}
(c)\quad At $\theta = 0$, $\cos\theta = 1$ and $\sec\theta = $ \underline{\hspace{1.5cm}}.
At $\theta = \pi/3$, $\cos\theta = 1/2$ and $\sec\theta = $ \underline{\hspace{1.5cm}}.
As the cosine value decreases toward 0, the secant value \underline{\hspace{2.5cm}}.

\vspace{0.12in}
\textbf{5. (AP-Style MC)}\quad Which of the following is NOT in the range of
$f(\theta) = \sec\theta$?

\begin{enumerate}[label=(\Alph*), leftmargin=2em, itemsep=3pt]
  \item $-3$
  \item $-1$
  \item $0$
  \item $1$
  \item $5$
\end{enumerate}

Answer: \underline{\hspace{1.5cm}}

\end{notesbox}

\vspace{0.08in}

\begin{notesbox}{Problem 6}

\textbf{6. (AP-Style MC)}\quad Which of the following gives the correct domain
restriction for $g(\theta) = \csc\theta$?

\begin{enumerate}[label=(\Alph*), leftmargin=2em, itemsep=3pt]
  \item $\theta \neq \pi/2 + n\pi$ for any integer $n$
  \item $\theta \neq n\pi$ for any integer $n$
  \item $\theta \neq \pi/4 + n\pi/2$ for any integer $n$
  \item $\theta \neq n\pi/2$ for any integer $n$
  \item There are no restrictions; $\csc\theta$ is defined for all real $\theta$
\end{enumerate}

Answer: \underline{\hspace{1.5cm}}

\end{notesbox}

\vspace{0.08in}

\begin{tcolorbox}[colback=navylight, colframe=navy, arc=2mm,
    left=3mm, right=3mm, top=2mm, bottom=2mm]
\small\textbf{Preview --- Lesson 3.12: Sinusoidal Function Contexts}\\
We now know all six trigonometric functions. In Lesson 3.12 we'll apply
sinusoidal models to real-world contexts: tides, seasonal temperature variation,
and periodic motion. Your fluency with the six functions --- including their
domains and ranges --- will be essential for interpreting and creating
these models.
\end{tcolorbox}

\end{document}
